\section{Introdução}
\label{sec:introducao}

\subsection{Contextualização e Problema de Pesquisa}

Ao longo das últimas décadas, o trabalho em escritório tem crescido de forma
significativa, impulsionado principalmente pelos avanços tecnológicos e pela
digitalização das atividades profissionais. Esse cenário foi intensificado após a
pandemia de COVID-19, que consolidou o modelo de trabalho remoto, permitindo que
diversas funções fossem executadas fora do ambiente corporativo tradicional,
muitas vezes apenas com o uso de um computador.

Entretanto, essa transformação trouxe novos desafios relacionados às condições de
trabalho, especialmente no que diz respeito à ergonomia. O ambiente doméstico, na
maioria dos casos, não é projetado para suportar longas jornadas de trabalho, o que
pode comprometer aspectos importantes como a postura, o conforto e a saúde do
trabalhador. Além disso, mesmo em escritórios convencionais, ainda são frequentes
as inadequações ergonômicas que contribuem para o surgimento de problemas de saúde
ocupacional.

Nesse contexto, os distúrbios osteomusculares relacionados ao trabalho (DORT)
destacam-se como um dos principais problemas de saúde entre trabalhadores de
escritório. Estudos indicam uma alta prevalência dessas condições, com índices que
podem variar entre 71,9\,\% e 80,81\,\% dos trabalhadores, afetando principalmente
regiões como pescoço, ombros e coluna lombar (OKEZUE et al., 2020; MOHAMMADIAN et
al., 2025). Esses distúrbios estão diretamente associados a fatores como postura
inadequada, permanência prolongada em posições estáticas, mobiliário inadequado e
ausência de pausas regulares.

Além dos fatores físicos, aspectos organizacionais e psicológicos, como o estresse
ocupacional, também desempenham um papel relevante no agravamento desses problemas.
A manutenção de posturas inadequadas por longos períodos pode reduzir a circulação
sanguínea, aumentar a tensão muscular e contribuir para o desenvolvimento de
condições crônicas, impactando diretamente a qualidade de vida e a produtividade
dos trabalhadores (MOHAMMADIAN et al., 2025).

Diante desse cenário, a ergonomia assume um papel fundamental na promoção da saúde
e segurança no ambiente de trabalho. Contudo, a adoção de práticas ergonômicas
adequadas ainda depende, em grande parte, da conscientização individual e de
intervenções pontuais, o que limita sua eficácia a longo prazo.

Nesse sentido, o uso de sistemas inteligentes surge como uma alternativa
promissora, permitindo o monitoramento contínuo da postura e a geração de feedback
em tempo real para o usuário. Tecnologias como sensores, dispositivos vestíveis e
algoritmos de análise de dados possibilitam identificar padrões de comportamento
inadequados e atuar de forma preventiva, contribuindo para a redução dos riscos
associados aos distúrbios osteomusculares.

Dessa forma, define-se como problema de pesquisa: \emph{como sistemas inteligentes
podem ser utilizados para monitorar e melhorar a ergonomia no ambiente de trabalho,
contribuindo para a prevenção de distúrbios osteomusculares?}

\subsection{Justificativa}

A crescente incidência de distúrbios osteomusculares entre trabalhadores evidencia
a necessidade de soluções eficazes voltadas à prevenção desses problemas. Além de
impactarem negativamente a saúde dos indivíduos, essas condições geram custos
significativos para empresas e sistemas de saúde, devido ao aumento do absenteísmo
e à redução da produtividade.

Com a expansão do trabalho remoto, a responsabilidade pela adequação ergonômica do
ambiente passou, em grande parte, para o próprio trabalhador, que nem sempre dispõe
de recursos ou conhecimento para implementar boas práticas posturais. Esse cenário
reforça a importância do desenvolvimento de ferramentas tecnológicas que auxiliem na
promoção da ergonomia de forma acessível e contínua.

Nesse contexto, os sistemas inteligentes aplicados à ergonomia apresentam grande
potencial, pois permitem o monitoramento individualizado e em tempo real,
promovendo a autocorreção postural e prevenindo o surgimento de lesões. Dessa forma,
este trabalho justifica-se pela necessidade de desenvolver soluções que contribuam
para a melhoria da saúde ocupacional, especialmente em ambientes de trabalho
sedentários.

\subsection{Objetivo Geral}

O objetivo deste trabalho é desenvolver um sistema inteligente capaz de identificar
posturas ergonomicamente inadequadas e fornecer alertas ao usuário em tempo real,
visando à prevenção de distúrbios osteomusculares.

\subsection{Objetivos Específicos}
\begin{itemize}[nosep]
  \item \textbf{(OE1)} Compreender as variações anatômicas dos usuários;
  \item \textbf{(OE2)} Desenvolvimento do sistema integrado com sensores para
        monitoramento postural;
  \item \textbf{(OE3)} Implementação do sistema de comunicação com o usuário;
  \item \textbf{(OE4)} Avaliar a eficiência do sistema proposto.
\end{itemize}

\noindent\textbf{Organização do trabalho.} O Capítulo~\ref{sec:revisao} revisa
a ergonomia de escritório, os DORT, os métodos de avaliação ergonômica, os
sensores usados em monitoramento postural e os trabalhos relacionados. O
Capítulo~\ref{sec:metodologia} descreve o sistema desenvolvido --- unidade
vestível, cadeira instrumentada, enlace sem fio, regra de alerta e alimentação
--- e os protocolos experimentais. O Capítulo~\ref{sec:resultados} apresenta a
caracterização dos sensores em bancada, as capturas do sistema completo com
múltiplos sujeitos, a avaliação da regra de fusão e a validação em hardware. O
Capítulo~\ref{sec:discussao} discute os resultados à luz dos objetivos
específicos e das limitações, e o Capítulo~\ref{sec:conclusao} conclui e aponta
trabalhos futuros.
